% ---------------------------------------------------------------------------
% Preambule -- PFE Toulouse INP-ENSEEIHT / MF2E
%
% Contraintes de forme imposees (student_context/Consignes pour les PFE2526.pdf) :
%   - coeur <= 30 pages, POLICE 12, figures comprises, hors annexes
%   - annexes <= 20 pages
%   - resume <= 2 pages en francais ET en anglais
%   - page de garde, remerciements, sommaire pagine, table des acronymes, biblio en fin
% Les pages de garde / biblio / sommaire / resume ne comptent pas dans les 30.
% ---------------------------------------------------------------------------

% --- Langue et encodage ----------------------------------------------------
\usepackage[T1]{fontenc}
\usepackage[utf8]{inputenc}
% `english` doit etre charge aussi : le resume anglais l'appelle via
% \begin{otherlanguage}{english}. Le DERNIER de la liste est la langue principale.
\usepackage[english,french]{babel}
\frenchsetup{SmallCapsFigTabCaptions=false}
\usepackage{csquotes}

% --- Fontes ----------------------------------------------------------------
% newtx (Times) : lisible en 12 pt et nettement plus compact que Latin Modern,
% ce qui compte quand 30 pages est une contrainte dure et non un objectif.
\usepackage{newtxtext}
\usepackage{newtxmath}
\usepackage{microtype}

% --- Mise en page ----------------------------------------------------------
% Marges laterales a 3 cm et non 2,5 : a 12 pt en Times, une justification de 16 cm donne
% ~95 signes par ligne, tres au-dela de la plage de confort (60-75). 15 cm ramene a ~85, ce
% qui reste long mais lisible, et le budget de pages interdit d'aller plus loin.
\usepackage[a4paper,top=2.5cm,bottom=2.5cm,left=3cm,right=3cm]{geometry}
\usepackage{setspace}
% 1,5 est une convention de BROUILLON annotable ; un document final se compose entre 1,15 et
% 1,25. Combine a la justification ci-dessus, le gain net est d'environ 20 % de texte par page
% -- ce qui n'est pas une facon de tricher avec les 30 pages mais de cesser d'en perdre.
\setstretch{1.15}

% Une veuve ou une orpheline est le defaut de composition le plus visible d'un rapport.
\widowpenalty=10000
\clubpenalty=10000
\raggedbottom
\usepackage{fancyhdr}
\pagestyle{fancy}
\fancyhf{}
\fancyhead[L]{\footnotesize\nouppercase{\leftmark}}
\fancyhead[R]{\footnotesize\thepage}
\renewcommand{\headrulewidth}{0.4pt}
\fancypagestyle{plain}{\fancyhf{}\fancyfoot[C]{\footnotesize\thepage}%
  \renewcommand{\headrulewidth}{0pt}}

% --- Tableaux et figures ---------------------------------------------------
\usepackage{graphicx}
\usepackage{booktabs}
\usepackage{tabularx}
\usepackage{multirow}
\usepackage{longtable}
\usepackage{array}
\usepackage[font=small,labelfont=bf,width=0.92\textwidth]{caption}
\usepackage{float}

% --- Mathematiques ---------------------------------------------------------
\usepackage{amsmath}
% newtxmath definit deja \Bbbk ; amssymb le redefinirait et LaTeX s'arrete.
\let\Bbbk\relax
\usepackage{amssymb}
\usepackage{mathtools}

% --- Nombres et unites -----------------------------------------------------
% Convention francaise : virgule decimale, espace insecable fine aux milliers.
\usepackage{siunitx}
\sisetup{
  locale = FR,
  group-separator = {\,},
  group-minimum-digits = 4,
  output-decimal-marker = {,},
  detect-all
}
\usepackage{textcomp}
% LE SYMBOLE EURO. Il apparait des centaines de fois dans ce memoire, il doit donc etre juste.
% newtx ne fournit pas d'euro en TS1 : \texteuro y declenche un "Faking \texteuro", et pdflatex
% fabrique alors le glyphe a la volee dans une fonte bitmap Computer Modern. Le resultat est
% visiblement etranger au texte -- graisse, chasse et dessin ne sont pas ceux de la page.
% `eurosym` fournit le glyphe officiel en Type 1, vectoriel, et s'accorde a un romain classique.
% NE PAS ecrire \DeclareSIUnit{\euro}{\text{\euro}} -- la definition serait recursive.
\usepackage{eurosym}
\newcommand{\eur}{\euro}
\DeclareSIUnit{\ha}{ha}
\DeclareSIUnit{\etp}{ETP}
\DeclareSIUnit{\tonne}{t}
% L'unite siunitx s'appelle \euros et NON \EUR : eurosym definit deja un \EUR{<montant>} qui
% prend un argument, et \DeclareSIUnit ne le remplace pas -- le symbole avalait alors
% l'accolade suivante ("Argument of \EUR has an extra }"). Le \text{} remet la macro en mode
% texte, siunitx composant ses unites en mode mathematique.
\DeclareSIUnit{\euros}{\text{\euro}}

% --- Listes ----------------------------------------------------------------
\usepackage{enumitem}
\setlist{noitemsep,topsep=4pt}

% --- Dessins ---------------------------------------------------------------
% tcolorbox charge deja pgf/tikz ; on l'expose explicitement pour les deux figures
% dessinees a la main (fil de justification, tornade de regret). Aucune dependance
% externe : pas de Graphviz, pas de pgfplots, pas de fichier de donnees a produire.
\usepackage{tikz}
\usetikzlibrary{arrows.meta,positioning}

% --- Palette des deux catalogues (fig:spectre, Ch. 4) ----------------------
% Un degrade « favorable -> defavorable » se dessine spontanement en VERT-JAUNE-ROUGE,
% qui est precisement la paire que 8 % des hommes ne distinguent pas (deuteranopie,
% protanopie). La palette ci-dessous garde la SEMANTIQUE (vert = favorable /
% agroecologique, jaune = temoin, rouge = defavorable / deregulateur) mais decale les
% deux extremes vers des teintes que le daltonisme rouge-vert separe encore : le vert
% tire vers le turquoise (composante bleue) et le rouge vers la brique (composante
% jaune), sur le modele de la rampe BrBG de ColorBrewer, declaree colorblind-safe.
% Les nuances intermediaires se composent par melange avec le sable (`scenVert!58!scenSable`),
% ce qui garantit une luminance croissante du centre vers les extremes -- second canal de
% lecture, independant de la teinte.
% La couleur n'est de toute facon JAMAIS le seul porteur d'information ici : les lignes
% sont ordonnees, et chacune porte son abreviation et son nom.
\definecolor{scenVert}{HTML}{2E9B7A}   % turquoise-vert : favorable / agroecologique
\definecolor{scenSable}{HTML}{F2E3B3}  % sable : le temoin, ni l'un ni l'autre
\definecolor{scenRouge}{HTML}{C4562F}  % brique : defavorable / deregulateur
\definecolor{scenGris}{HTML}{E2E2E2}   % hors gradient (P9)
% Couleurs de cadre, prises a la palette Okabe-Ito (concue pour le daltonisme). Elles
% sont DOUBLEES d'un style de trait (plein / tirete / pointille) pour que les groupes
% restent distincts meme imprimes en noir et blanc.
\definecolor{cadreClimat}{HTML}{0072B2} % bleu
\definecolor{cadreMarche}{HTML}{6A3D9A} % violet

% --- Boites ----------------------------------------------------------------
\usepackage[most]{tcolorbox}

% --- MINI-SOMMAIRE : l'entree en matiere de chaque chapitre -----------------
% Un filet en haut, un filet en bas, une seule ligne courante : le lecteur sait ou il va
% sans qu'on lui vende une demi-page de table des matieres locale. Le budget des 30 pages
% interdit une liste a puces ici -- sept chapitres x une demi-page couteraient 3 pages.
% Usage :
%   \minisommaire{\ms{sec:dispositif}{ce qui compose un scenario} \ms{sec:grille}{...}}
% \ms compose le numero de section en gras puis son resume ; le separateur est automatique.
\newtcolorbox{minisommaireboite}{%
  enhanced, sharp corners, breakable,
  colback = black!3, colframe = black!45,
  boxrule = 0pt, toprule = 0.7pt, bottomrule = 0.7pt,
  left = 6pt, right = 6pt, top = 4pt, bottom = 4pt,
  fontupper = \footnotesize,
}
\newif\ifmspremier
\newcommand{\ms}[2]{%
  \ifmspremier\mspremierfalse\else\ \textperiodcentered\ \fi
  \textbf{\ref{#1}}~#2}
\newcommand{\minisommaire}[1]{%
  \begin{minisommaireboite}%
    \mspremiertrue
    {\scriptsize\sffamily\bfseries DANS CE CHAPITRE}\quad #1%
  \end{minisommaireboite}%
  \par\smallskip}

% --- Bibliographie ---------------------------------------------------------
\usepackage[backend=biber,style=authoryear,maxcitenames=2,maxbibnames=99,
            giveninits=true,uniquename=init,sorting=nyt]{biblatex}
\addbibresource{biblio.bib}

% --- Liens -----------------------------------------------------------------
\usepackage[hidelinks,breaklinks]{hyperref}
% Metadonnees du PDF : ce qui s'affiche dans la barre de titre du lecteur et dans les
% proprietes du fichier. Un PDF sans metadonnees s'appelle "memoire.pdf" partout.
\hypersetup{
  pdftitle   = {Porter, calibrer et explorer un modele d'allocation de cultures --
                MOSAICA en Python/Pyomo (Guadeloupe)},
  pdfauthor  = {Clement Le Borgne-Lariviere},
  pdfsubject = {Projet de Fin d'Etudes -- Toulouse INP-ENSEEIHT, departement MF2E},
  pdfkeywords= {MOSAICA; allocation de cultures; MILP; Pyomo; HiGHS; calibration;
                prospective; Guadeloupe; agroecologie},
  pdfcreator = {pdflatex}
}
% CLEVEREF A ETE RETIRE (2026-08-15). Dans cette installation, `\cref` sur une reference
% AVANT-COUREUSE (label defini plus loin dans le document) plante avec "Missing \endcsname
% inserted" des la premiere passe -- donc fatalement, puisque la premiere passe est celle qui
% ecrit le .aux. Le document ne peut jamais converger. Reproduit sur un fichier minimal :
% `\cref` sur un label DEJA defini passe, sur un label a venir non.
% La convention du memoire est donc explicite et francaise, ce qui est de toute facon plus
% lisible qu'un \cref sous babel-french :
%     chapitre~\ref{ch:xxx}   §~\ref{sec:xxx}   tableau~\ref{tab:xxx}
%     figure~\ref{fig:xxx}    annexe~\ref{ann:xxx}   equation~\eqref{eq:xxx}
% NE PAS remettre \usepackage{cleveref} sans avoir reproduit le test ci-dessus.

% --- Acronymes -------------------------------------------------------------
\usepackage[acronym,nomain,nonumberlist,nopostdot]{glossaries}
\makenoidxglossaries

% --- Titres de chapitre sobres --------------------------------------------
\usepackage{titlesec}
% CHAQUE CHAPITRE COMMENCE SUR UNE PAGE NEUVE. La classe titlesec est `top` (titre en haut
% d'une page neuve, comportement de \chapter) et NON `page`, qui reserverait une page entiere
% au seul titre comme le fait \part -- et qui dedouble les entrees de sommaire des chapitres
% non numerotes. Le titre reste compose sobrement, sans le pave de blanc de 2,5 cm que `report`
% place au-dessus par defaut.
\titleclass{\chapter}{top}
\titleformat{\chapter}[block]
  {\normalfont\Large\bfseries}{\normalsize\mdseries\textsc{\chaptertitlename~\thechapter}\\[2pt]}{0pt}{\Large}
\titlespacing*{\chapter}{0pt}{0pt}{16pt}
\titlespacing*{\section}{0pt}{14pt plus 3pt minus 2pt}{6pt}

% --- Raccourcis ------------------------------------------------------------
\newcommand{\mosaica}{\textsc{Mosaica}}
\newcommand{\gams}{\textsc{Gams}}
\newcommand{\highs}{\textsc{HiGHS}}
\newcommand{\pad}{\textsc{pad}}

% Compte de pages du coeur : le compteur arabe redemarre a 1 au chapitre 1 et
% \pageref{fin-coeur} donne donc directement le nombre de pages du coeur. Il est
% ecrit dans le .log a chaque compilation pour que la contrainte des 30 pages
% soit verifiee mecaniquement et non a l'oeil.
% `\typeout` n'expanse PAS `\pageref` -- il ecrivait la macro litteralement, et le
% controle des 30 pages ne controlait donc rien (corrige le 2026-08-11). `refcount`
% fournit la version developpable, seule utilisable dans un \typeout.
\usepackage{refcount}
\newcommand{\annoncecoeur}{%
  \typeout{^^J===== PAGES DU COEUR : \getpagerefnumber{fin-coeur} / 30 =====^^J}}

